\chapter{Networking for Trading}\label{ch:nw:networking-for-trading}

The operator of the consolidated U.S. options feed tells its subscribers how much network to buy, and it does not quote an
average. For July 2026 it projected 4.403 gigabits in the busiest 100 milliseconds and 0.501 gigabits in the busiest 10
milliseconds, per copy of the feed: \qty{44}{\giga\bit\per\second} over the longer window, \qty{50}{\giga\bit\per\second}
over the shorter, twice that for a subscriber that takes both redundant copies. Its own statistics say the same thing from the
other side: the busiest millisecond of July 2026 carried messages at five times the rate of the busiest second. A network
sized for the average of a trading day loses packets in its busiest milliseconds, and those are the milliseconds a trading firm
is paid to see.

This chapter opens the book's first part, the network between a venue and the firm's servers. It prices a byte on the wire,
shows how one packet reaches many subscribers, explains why market data travels in datagrams and orders in streams, and ends
where most losses begin: at a switch port that receives more than it can send.

\section{The Ethernet frame and what a byte costs on the wire}

\begin{definition}[Network interface card, Ethernet frame]\label{def:nw:networking-for-trading:frame}
A \emph{network interface card}\index{network interface card} (NIC) is the device that connects a computer to a network link:
it turns the bits on the wire into data in the host's memory and back. An \emph{Ethernet frame}\index{Ethernet frame} is the
unit an Ethernet link carries: a destination and a source address (6 bytes each), a type (2 bytes), a payload of 46 to 1\,500
bytes and a frame check sequence (4 bytes), from 64 to 1\,518 bytes in all without a VLAN tag. On the wire each frame is
preceded by 8 bytes of preamble and start delimiter and followed by at least 12 bytes of idle inter-frame gap.
\end{definition}

\begin{definition}[Line rate, serialisation delay, maximum transmission unit]\label{def:nw:networking-for-trading:ser}
The \emph{line rate}\index{line rate} $R_{\mathrm L}$ of a link is the rate at which it clocks bits onto the medium (1, 10, 25
or \qty{100}{\giga\bit\per\second}). The \emph{serialisation delay}\index{serialisation delay} of a frame is the time the link
takes to put it on the wire. The \emph{maximum transmission unit}\index{maximum transmission unit} (MTU) is the largest payload
a link carries in one frame, 1\,500 bytes for standard Ethernet.
\end{definition}

\begin{proposition}[What a frame costs]\label{prop:nw:networking-for-trading:ser}
A frame of $\ell_{\mathrm f}$ bytes occupies $\ell_{\mathrm f} + 20$ bytes of line time, so its \omterm{def:nw:networking-for-trading:ser}{serialisation delay} is
\[
  t_{\mathrm{ser}} = \frac{8\,(\ell_{\mathrm f} + 20)}{R_{\mathrm L}},
\]
and a link carries at most $R_{\mathrm L} / (8(\ell_{\mathrm f}+20))$ frames a second. Every device that must receive a whole
frame before it sends it on pays $t_{\mathrm{ser}}$ once more.
\end{proposition}

\begin{proof}
The preamble, start delimiter and inter-frame gap are 20 bytes of line time that belong to each frame and to no other; a link
sends one bit every $1/R_{\mathrm L}$. A store-and-forward device (chapter~2) starts to transmit only once the last bit has
arrived, a delay of one serialisation at the incoming rate.
\end{proof}

\begin{omfigure}
\begin{tikzpicture}[font=\scriptsize,f/.style={draw,minimum height=0.95cm,minimum width=1.32cm,text width=1.2cm,
  align=center,inner sep=1pt}]
\foreach \i/\n/\b/\c in {0/{preamble\\+ SFD}/8/gray!15, 1/Ethernet\\header/14/omDef!15, 2/IPv4\\header/20/omDef!15,
  3/UDP\\header/8/omDef!15, 4/{Mold-\\UDP64}/20/omMeth!18, 5/length/2/omMeth!18, 6/add\\order/36/omThm!25,
  7/FCS/4/omDef!15, 8/{gap}/12/gray!15}
  {\node[f,fill=\c,anchor=south west] at (1.35*\i,0) {\n}; \node[font=\scriptsize] at (1.35*\i+0.66,-0.25) {\b\ bytes};}
\draw[decorate,decoration={brace,amplitude=4pt,mirror}] (1.35,-0.5) -- (10.8,-0.5)
  node[midway,below=4pt,font=\scriptsize] {Ethernet frame: 104 bytes (64 minimum, 1\,518 maximum)};
\draw[decorate,decoration={brace,amplitude=4pt}] (0,1.05) -- (12.15,1.05)
  node[midway,above=4pt,font=\scriptsize] {line time: 124 bytes $=$ \qty{99.2}{\nano\second} at \qty{10}{\giga\bit\per\second},
  \qty{39.7}{\nano\second} at \qty{25}{\giga\bit\per\second}};
\end{tikzpicture}

\omcaption{One add-order message on the wire, in the simulator's feed format (a MoldUDP64 packet in a UDP datagram in an
\omterm{def:nw:networking-for-trading:frame}{Ethernet frame}). The 36 bytes the strategy reads are 29\% of the line time the frame occupies. Boxes are not to scale; byte
counts from \cref{prop:nw:networking-for-trading:ser} and the formats' headers.}\label{fig:nw:networking-for-trading:frame}
\end{omfigure}

\begin{example}[One add order]\label{ex:nw:networking-for-trading:mold}
The exchange simulator of One Quant Book~10 publishes its feed as MoldUDP64 packets: a 20-byte header (session, sequence
number, message count), then each message behind a 2-byte length. An add-order message is 36 bytes, so a packet carrying one
is 58 bytes of UDP payload; with 8 bytes of UDP header, 20 of IPv4 header and 18 of Ethernet header and check sequence the
frame is 104 bytes and occupies 124 bytes of line time (\cref{fig:nw:networking-for-trading:frame}): \qty{99.2}{\nano\second}
at \qty{10}{\giga\bit\per\second}, \qty{39.7}{\nano\second} at 25. A venue that packs eight such messages in one packet
spends the 66 bytes of fixed overhead once for eight messages instead of eight times.
\end{example}

\begin{table}[ht]
\centering\small
\begin{tabular}{rrrrr}
\toprule
Frame (bytes) & \qty{1}{\giga\bit\per\second} & \qty{10}{\giga\bit\per\second} & \qty{25}{\giga\bit\per\second} &
\qty{100}{\giga\bit\per\second} \\
\midrule
64 & 672.0 & 67.2 & 26.9 & 6.7 \\
128 & 1\,184.0 & 118.4 & 47.4 & 11.8 \\
256 & 2\,208.0 & 220.8 & 88.3 & 22.1 \\
512 & 4\,256.0 & 425.6 & 170.2 & 42.6 \\
1\,024 & 8\,352.0 & 835.2 & 334.1 & 83.5 \\
1\,518 & 12\,304.0 & 1\,230.4 & 492.2 & 123.0 \\
\bottomrule
\end{tabular}
\caption{\omterm{def:nw:networking-for-trading:ser}{Serialisation delay} in nanoseconds by frame size and \omterm{def:nw:networking-for-trading:ser}{line rate} (\cref{prop:nw:networking-for-trading:ser}, preamble
and inter-frame gap included). Computed by \texttt{nw\_wire.frame\_table}.}\label{tab:nw:networking-for-trading:ser}
\end{table}

\Cref{tab:nw:networking-for-trading:ser} is the first latency table of this book and the one most often forgotten. A
full-size frame at \qty{10}{\giga\bit\per\second} takes \qty{1.23}{\micro\second} to serialise, half of the
\qty{2.5}{\micro\second} that an engineer of a large market maker gave in a public talk as a very good wire-to-wire time for a
whole software trading system (One Quant Book~13, chapter~1). Moving from 10 to \qty{25}{\giga\bit\per\second} saves 60\% of
every serialisation, a reason to buy the faster port even when the slower one would carry the traffic.
Serialisation is also why small frames matter: an order is rarely larger than a hundred bytes, and a switch that forwards it
only once it has all of it pays for all of it.

\begin{remark}[Bits in flight]\label{rem:nw:networking-for-trading:flight}
Light in standard single-mode fibre travels a metre in about \qty{4.88}{\nano\second} (group index 1.462, chapter~10).
A 1\,518-byte frame at \qty{10}{\giga\bit\per\second} is therefore 252~metres long in the fibre, longer than any cable in a
data-centre hall: its first bit reaches the far end long before its last bit has left. Distance and serialisation are separate
terms of a latency budget, and at the scale of a building serialisation usually wins.
\end{remark}

\omcode{code/firm/netsim/firm_netsim.py}{36}{50}{Frame arithmetic in \texttt{firm.netsim}: the smallest frame is padded to 64
bytes, and every frame pays 20 bytes of preamble and gap.}\label{lst:nw:networking-for-trading:frame}

\section{IP multicast: groups, membership and snooping}

A venue sends the same market data to hundreds of subscribers. Sending one copy to each over its own connection would make
the last subscriber wait for all the others (chapter~12 studies a venue that did exactly that); sending one copy that the
network replicates makes them all equal to within the replication.

\begin{definition}[Multicast group, IGMP]\label{def:nw:networking-for-trading:group}
A \emph{multicast group}\index{multicast group} is an IPv4 address between 224.0.0.0 and 239.255.255.255 that names a set of
receivers rather than one host: a datagram sent to it is delivered to every member. The \emph{Internet Group Management
Protocol}\index{Internet Group Management Protocol} (IGMP) is the protocol by which a host tells its neighbouring routers
which groups it wants to receive, and, in version~3, from which sources.
\end{definition}

A venue's feed is usually spread over many groups, one per channel or partition of instruments, so that a subscriber
receives only the groups it needs (chapter~22 sizes the options feeds this way). A group must also be delivered on the
Ethernet segment, and Ethernet addresses are 48 bits: a group address is mapped to the Ethernet address
\texttt{01:00:5e} followed by the low 23 bits of the group, so the upper 5 of the 28 bits that distinguish groups are lost and
32 groups share each Ethernet address. A card filtering on Ethernet addresses may therefore let through groups the host never
joined, and the kernel discards them after having paid for them.

\begin{definition}[IGMP snooping]\label{def:nw:networking-for-trading:snooping}
\emph{IGMP snooping}\index{IGMP snooping} is a switch's practice of reading the IGMP membership reports that cross it and
forwarding each group's traffic only to the ports on which a member has reported, instead of flooding every multicast frame
to every port.
\end{definition}

\begin{omfigure}
\begin{tikzpicture}[font=\footnotesize,b/.style={draw,rounded corners=2pt,minimum height=0.65cm,align=center,inner sep=3pt},
  >=Stealth]
\node[b,fill=gray!10] (v) at (0,0) {venue feed\\ groups A, B};
\node[b,fill=omDef!15,minimum width=1.9cm,minimum height=3.6cm] (s) at (3.6,0) {switch\\ (IGMP\\ snooping)};
\foreach \y/\t/\g in {1.35/{server 1: joined A}/{A}, 0.45/{server 2: joined A and B}/{A, B}, -0.45/{server 3: joined B}/{B},
  -1.35/{server 4: joined nothing}/{nothing}}
  {\node[b,minimum width=4.1cm,anchor=west] (h) at (7.4,\y) {\t};
   \draw[->,thick] (s.east|-h.west) -- node[above,font=\scriptsize]{sends \g} (h.west);}
\draw[->,thick] (v) -- node[above,font=\scriptsize]{A and B} (s.west|-v.east);
\end{tikzpicture}

\omcaption{\omterm{def:nw:networking-for-trading:snooping}{IGMP snooping}. The switch learns from the servers' membership reports which ports want which group and replicates
each group only to them; server~4 receives nothing. Without snooping, all four ports would carry both groups, and servers that
joined nothing would still spend their cards' and kernels' time discarding them.}\label{fig:nw:networking-for-trading:snoop}
\end{omfigure}

\begin{example}[Two groups, one Ethernet address]\label{ex:nw:networking-for-trading:mac}
Group 239.255.1.2 maps to \texttt{01:00:5e:7f:01:02}: the low 23 bits of 255.1.2 are 127.1.2 once the top bit of 255 is
dropped. Group 224.127.1.2 maps to the same address. A server that joined only the first receives, from any switch or card
that filters by Ethernet address alone, the traffic of the second as well. Venues allocate their groups with this in mind, and
a firm that runs its own internal multicast (chapter~27's ticker plant) should too.
\end{example}

\section{Datagrams against streams}

Two transport protocols carry almost everything in a trading network. UDP sends datagrams: each is delivered whole or not at
all, with no ordering and no retransmission, and one datagram sent to a group reaches every member. TCP carries a byte
stream between two endpoints: it numbers the bytes, acknowledges them, retransmits what is lost and delivers everything in
order.

Market data travels in UDP multicast because only multicast sends one copy to all, and because a feed must not slow down
for its slowest subscriber: TCP's flow control would do exactly that. The price is that loss is the application's problem. The
venue sends every packet on two independent lines (One Quant Book~10, chapter~26), numbers them, and runs retransmission and
snapshot services; the firm's handler arbitrates and recovers (One Quant Book~13, chapters~16 and~18). Orders travel over TCP
because an order must arrive exactly once, in order with the firm's other orders, and because the venue must know which
connection sent it: the session layer of the order-entry protocol builds on those guarantees.

\begin{definition}[Head-of-line blocking]\label{def:nw:networking-for-trading:hol}
\emph{Head-of-line blocking}\index{head-of-line blocking} is the delay imposed on every item in an ordered queue by one item at
its head that cannot proceed. In TCP, a lost segment holds back every later byte of the stream, already received or not, until
its retransmission arrives.
\end{definition}

\omterm{def:nw:networking-for-trading:hol}{Head-of-line blocking} is the cost of TCP's ordering. A cancel sent a microsecond after a lost new-order message waits for the
new order's retransmission, which comes after at least a round trip and usually after a timer: on a clean colocation link loss
is rare, and when it happens the whole session stalls. Ordering is also a feature: the venue must never see the cancel before
the order it cancels.

\begin{definition}[Nagle's algorithm]\label{def:nw:networking-for-trading:nagle}
\emph{Nagle's algorithm}\index{Nagle's algorithm} is TCP's rule that a sender holds new data back, instead of sending it in a
small segment, while any data it has already sent remains unacknowledged. It saves bandwidth when a program writes many small
pieces, and was designed for interactive terminals.
\end{definition}

\begin{proposition}[Nagle meets the delayed acknowledgement]\label{prop:nw:networking-for-trading:nagle}
Let a sender write a small order and, before it is acknowledged, a second small order, on a connection with \omterm{def:nw:networking-for-trading:nagle}{Nagle's algorithm}
enabled, to a receiver that delays its acknowledgements (TCP allows up to half a second). The second order leaves only when
the first is acknowledged: after the receiver's delayed-acknowledgement timer, or a round trip if the receiver has data of its
own to send. Disabling \omterm{def:nw:networking-for-trading:nagle}{Nagle's algorithm} (the socket option \texttt{TCP\_NODELAY}) removes the wait.
\end{proposition}

\begin{proof}
By the rule, the second write is held while the first segment is unacknowledged. The receiver sends its acknowledgement when
its timer fires, when it has a second full segment to acknowledge (it has not: one small segment arrived), or with its own data.
\end{proof}

\begin{method}[An order-entry socket]\label{met:nw:networking-for-trading:socket}
\begin{enumerate}
\item Set \texttt{TCP\_NODELAY} on every order-entry and drop-copy connection.
\item Write each message with one call, so that one message is one segment and nothing waits for a later write.
\item Where the venue's session protocol acknowledges at the application level, acknowledge promptly at the TCP level too
  (Linux's \texttt{TCP\_QUICKACK} is not permanent and must be set again after reads).
\item Keep the connection up (the session layer's heartbeats) and pre-warmed: a new TCP connection costs a round trip before
  its first byte, and chapter~20 adds a TLS handshake on top for crypto venues.
\end{enumerate}
\end{method}

\section{Congestion: egress queues, incast and microbursts}

A switch forwards each frame to an output port. If frames destined to one port arrive faster than the port's \omterm{def:nw:networking-for-trading:ser}{line rate}, they
wait.

\begin{definition}[Egress queue, oversubscription ratio]\label{def:nw:networking-for-trading:egress}
An \emph{egress queue}\index{egress queue} is the buffer in which a switch holds frames waiting to be sent on an output port;
when it is full, arriving frames are dropped (tail drop). The \emph{oversubscription ratio}\index{oversubscription ratio} of a
port is the sum of the \omterm{def:nw:networking-for-trading:ser}{line rates} that can send to it divided by its own \omterm{def:nw:networking-for-trading:ser}{line rate}.
\end{definition}

\begin{definition}[Incast, microburst]\label{def:nw:networking-for-trading:burst}
\emph{Incast}\index{incast} is the congestion that arises when many senders transmit to one receiver at the same moment. A
\emph{microburst}\index{microburst} is a period of microseconds to milliseconds during which the traffic offered to a port
exceeds its \omterm{def:nw:networking-for-trading:ser}{line rate}, while its average over any interval a monitoring tool reports (a second, a minute) stays well below
it.
\end{definition}

A feed that uses half of a port on average can overflow it for a millisecond, and a trading network is built to make this
happen: every venue's feed, and every partition of every feed, reacts to the same market events at the same time.

\begin{proposition}[The fluid bound]\label{prop:nw:networking-for-trading:fluid}
Let $n$ inputs send at rates $r_1, \dots, r_n$ to a port of \omterm{def:nw:networking-for-trading:ser}{line rate} $R_{\mathrm L}$ during a burst of length $D$, starting
from an empty queue. At the end of the burst the queue holds
\[
  B = \Bigl(\textstyle\sum_i r_i - R_{\mathrm L}\Bigr)^{+} D / 8 \text{ bytes},
\]
and the last byte of the burst leaves $8B/R_{\mathrm L}$ after it arrived. Tail drop loses nothing if and only if the buffer
available to the port is at least $B$; otherwise it loses $B$ minus the buffer, in whole frames.
\end{proposition}

\begin{proof}
During the burst the queue grows at the difference of the rates when it is positive and stays empty otherwise; the port then
drains it at $R_{\mathrm L}$. The discrete version, frame by frame, is the loop of \cref{lst:nw:networking-for-trading:egress},
which reduces to the fluid formula when frames are small against $B$.
\end{proof}

\begin{omfigure}
\begin{tikzpicture}
\begin{axis}[width=11.5cm,height=5.4cm,xmode=log,xmin=0.7,xmax=1400,ymin=0,ymax=380,
  xtick={1,10,100,1000},xticklabels={1~ms,10~ms,100~ms,1~s},
  xlabel={averaging window (log scale)},ylabel={peak rate (million messages / s)},
  tick label style={font=\small},label style={font=\small},grid=major,grid style={gray!25},
  legend columns=3,legend style={font=\footnotesize,at={(0.5,-0.32)},anchor=north,draw=none,fill=none,
  /tikz/every even column/.append style={column sep=8pt}},table/col sep=comma]
\addplot[omExo,very thick,mark=*,mark size=1.6pt] table[x=window_ms,y=2024_01]{figdata/networks/01-networking-for-trading/opra.csv};
\addplot[omMeth,very thick,mark=square*,mark size=1.6pt] table[x=window_ms,y=2025_03]{figdata/networks/01-networking-for-trading/opra.csv};
\addplot[omThm,very thick,mark=triangle*,mark size=2pt] table[x=window_ms,y=2026_07]{figdata/networks/01-networking-for-trading/opra.csv};
\legend{January 2024,March 2025,July 2026}
\end{axis}
\end{tikzpicture}

\omcaption{The consolidated U.S. options feed's busiest window of each length, as a rate: the peak count of messages in one
second, 100~ms, 10~ms and 1~ms divided by the window. The shorter the window, the higher the rate; in July 2026 the busiest
millisecond ran at 5.5 times the busiest second. Data: OPRA, \emph{Key Operating Metrics} (August 2026), transcribed in
\texttt{data/networks/opra\_metrics.csv}.}\label{fig:nw:networking-for-trading:opra}
\end{omfigure}

\Cref{fig:nw:networking-for-trading:opra} is the published version of the \omterm{def:nw:networking-for-trading:burst}{microburst}: the same feed, in the same month, is
63.9 million messages a second or 350 million a second depending on the window one reads it through. The operator's capacity
notice, which subscribers use to size their links, is written in 10-millisecond peaks for that reason
(\cref{dat:nw:networking-for-trading:opra}); a switch port does not average over 10 milliseconds, it queues frame by frame.

\begin{dated}{2026-09}{How much network the options feed needs}\label{dat:nw:networking-for-trading:opra}
SIAC's capacity projection for OPRA effective July 2026 gives, for one copy of the feed, 13.575 million messages, 4.403
gigabits and 1.564 million packets in the busiest 100 milliseconds, and 1.562 million messages, 0.501 gigabits and 169\,000
packets in the busiest 10 milliseconds. That is \qty{44.0}{\giga\bit\per\second} and \qty{50.1}{\giga\bit\per\second}; about 352
bytes and 8.7 messages per packet over the 100-millisecond window. Subscribers taking both redundant copies double the
bandwidth, and the notice asks for 10\% more for retransmissions. The notice states the 10-millisecond interval ``reflects system
utilization during bursts of traffic''.
\end{dated}

To see the mechanism rather than its footprint, the chapter's tutorial simulates eight feeds into one
\qty{10}{\giga\bit\per\second} port with \texttt{firm.netsim}. Each feed averages 400\,000 packets a second in clusters (a
market event produces a run of messages), one to thirty add-order messages a packet; the eight together offer about
\qty{5.2}{\giga\bit\per\second}, half the port. In one version every feed's clusters are independent. In the other, 30\% of
each feed's packets come in clusters triggered by common events, 400 a second, that hit all eight feeds at once, as a macro
release or a large trade in an index future does.

\begin{omfigure}
\begin{tikzpicture}
\begin{groupplot}[group style={group size=2 by 1,horizontal sep=1.7cm},width=6.6cm,height=5.2cm,
  tick label style={font=\small},label style={font=\small},grid=major,grid style={gray!25},table/col sep=comma]
\nextgroupplot[xmode=log,ymode=log,xmin=0.7,xmax=14000,ymin=3,ymax=300,xlabel={window (\unit{\micro\second}, log)},
  ylabel={peak rate (\unit{\giga\bit\per\second}, log)},xtick={1,10,100,1000,10000},
  xticklabels={1,10,100,1\,000,10\,000},ytick={5,10,20,50,100,200},yticklabels={5,10,20,50,100,200}]
\addplot[omDef,very thick,mark=*,mark size=1.4pt] table[x=window_us,y=independent]{figdata/networks/01-networking-for-trading/peak.csv};
\addplot[omThm,very thick,mark=square*,mark size=1.4pt] table[x=window_us,y=common]{figdata/networks/01-networking-for-trading/peak.csv};
\addplot[gray,dashed,domain=0.7:14000] {10};
\nextgroupplot[xmode=log,xmin=8,xmax=6000,ymin=0,ymax=26,xlabel={egress buffer (KB, log)},ylabel={frames lost (\%)},
  xtick={10,100,1000},xticklabels={10,100,1\,000},
  legend columns=2,legend style={font=\footnotesize,at={(-0.2,-0.32)},anchor=north,draw=none,fill=none,
  /tikz/every even column/.append style={column sep=8pt}}]
\addplot[omDef,thick,dashed,mark=o,mark size=1.4pt] table[x=buffer_kb,y=f4_independent]{figdata/networks/01-networking-for-trading/fanin.csv};
\addplot[omThm,thick,dashed,mark=square,mark size=1.4pt] table[x=buffer_kb,y=f4_common]{figdata/networks/01-networking-for-trading/fanin.csv};
\addplot[omDef,very thick,mark=*,mark size=1.4pt] table[x=buffer_kb,y=f8_independent]{figdata/networks/01-networking-for-trading/fanin.csv};
\addplot[omThm,very thick,mark=square*,mark size=1.4pt] table[x=buffer_kb,y=f8_common]{figdata/networks/01-networking-for-trading/fanin.csv};
\legend{4 feeds independent,4 feeds with common events,8 feeds independent,8 feeds with common events}
\end{groupplot}
\end{tikzpicture}

\omcaption{Simulation: feeds into one \qty{10}{\giga\bit\per\second} egress port. Left: the busiest window of each length for eight
feeds (dashed: the \omterm{def:nw:networking-for-trading:ser}{line rate}); independent bursts stay under it at 100~\unit{\micro\second} and beyond, common events push it
past \qty{60}{\giga\bit\per\second}. Right: frames lost to tail drop by buffer size, 50~ms of traffic. The same average load
needs about 30~KB of buffer when bursts are independent and about a megabyte when they are common. Data:
\texttt{fig\_bursts.py} on \texttt{firm.netsim} (seeded, deterministic).}\label{fig:nw:networking-for-trading:fanin}
\end{omfigure}

The two versions have the same average load, and the difference is all in the correlation
(\cref{fig:nw:networking-for-trading:fanin}). With independent bursts the busiest 100 microseconds carry
\qty{9.1}{\giga\bit\per\second}, and the buffer that would have lost nothing, the largest backlog of the run, is 21 to 31~KB
over twenty seeds. With common events the busiest 100 microseconds carry \qty{68}{\giga\bit\per\second}, and that buffer is
between 0.44 and 2.4~MB over the same twenty seeds, 1.1~MB on average. Removing the common events is the ablation: the
megabyte is caused by the correlation, not by the load. A port's traffic statistics averaged over a second show half a link in
both cases.

\begin{dated}{2026-09}{A low-latency switch's buffer}\label{dat:nw:networking-for-trading:buffer}
Cisco's data sheet for the Nexus 3548-X, a 48-port \qty{10}{\giga\bit\per\second} switch sold for low latency, gives its
buffer as ``6 MB shared among 16 ports; 18 MB total'' and latencies ``as low as'' \qty{250}{\nano\second} in normal mode and
\qty{200}{\nano\second} in its warp mode. A port whose fifteen neighbours are also busy can count on about a sixteenth of its
6~MB block, 375~KB.
\end{dated}

Read against the simulation, the box says that a port which can borrow the whole shared block survives the common-event
bursts of \cref{fig:nw:networking-for-trading:fanin}, and a port that gets only its sixteenth does not: at 500~KB of buffer the
eight feeds with common events still lose 2.6\% of their frames. The events that fill every port at once are the events that
leave each port with only its share. The remedies are a faster egress port (the eight feeds' common bursts reach
\qty{16}{\giga\bit\per\second} over a millisecond, within a \qty{25}{\giga\bit\per\second} port), fewer feeds per port, or a
switch whose buffer is dedicated rather than shared; chapter~2 adds the devices that do not queue at all.

\section{Tutorial: microbursts from frames to drops}\label{tut:nw:networking-for-trading:tut}

\begin{tutorial}
\textbf{Goal.} Price frames on the wire, read the published burst curve, and simulate a fan-in port until it drops.
\textbf{End state:} \cref{tab:nw:networking-for-trading:ser}, \cref{fig:nw:networking-for-trading:opra,fig:nw:networking-for-trading:fanin},
and the zero-loss buffers quoted in the text.
\begin{tutsteps}
\item \textbf{Frames.} \texttt{nw\_wire.frame\_table()} and \texttt{mold\_example()} reproduce the table and
  \cref{ex:nw:networking-for-trading:mold} from \cref{lst:nw:networking-for-trading:frame}.
\item \textbf{The published curve.} \texttt{opra\_curve()} turns the operator's peak counts into rates; \texttt{opra\_capacity()}
  turns the capacity notice into gigabits a second and bytes a packet.
\item \textbf{Traffic.} \texttt{feed()} draws one feed's packets from \texttt{firm.netsim.bursty\_arrivals}, and
  \texttt{fanin\_traffic()} merges eight of them, with or without common events (\cref{lst:nw:networking-for-trading:feed}).
\item \textbf{The port.} \texttt{firm.netsim.egress} runs the tail-drop queue (\cref{lst:nw:networking-for-trading:egress});
  \texttt{fanin()} sweeps the buffer, \texttt{zero\_loss\_buffer()} runs twenty seeds with an unbounded buffer and returns each
  run's largest backlog. \texttt{python fig\_bursts.py} writes the charts' data.
\end{tutsteps}
\textbf{What to change next.} Vary the share of packets that come with common events from 0 to 0.5 and plot the zero-loss
buffer against it; replace the \qty{10}{\giga\bit\per\second} port by a \qty{25}{\giga\bit\per\second} one and find the smallest
buffer that loses nothing.
\end{tutorial}

\omcode{code/networks/01-networking-for-trading/python/nw_wire.py}{62}{77}{One feed and a fan-in of eight: the common event
times are drawn once and shared by every feed.}\label{lst:nw:networking-for-trading:feed}

\omcode{code/firm/netsim/firm_netsim.py}{119}{140}{The egress queue: the backlog drains at the line rate between arrivals, and a
frame that does not fit is dropped whole.}\label{lst:nw:networking-for-trading:egress}

\section{Build: the network model}\label{bld:nw:networking-for-trading:netsim}

\begin{build}
\bfield{Purpose} A packet-level model of the firm's network that every later chapter of Parts~I and~II uses: frames, links,
queues, switches and multicast, with exact byte accounting and seeded traffic.
\bfield{Interface} \texttt{firm\_netsim}: constants for every header; \texttt{frame\_bytes}, \texttt{mold\_payload},
\texttt{wire\_bytes}, \texttt{ser\_ns}, \texttt{prop\_ns}; \texttt{mcast\_mac}; \texttt{peak\_rate(t, bytes, window)};
\texttt{event\_times}, \texttt{bursty\_arrivals}; \texttt{egress(t, frames, gbps, buffer)} returning departures, drops and
backlogs; \texttt{forward\_ns(frame, gbps, mode, fabric\_ns)} for store-and-forward, cut-through and layer-1 devices;
\texttt{replicate(groups, members, ports, snooping)}; \texttt{merge}. C++20 twin of the \omterm{def:nw:networking-for-trading:egress}{egress queue},
\texttt{cpp/firm\_netsim.hpp}.
\bfield{Rules} Integer nanoseconds and bytes; every frame pays 20 bytes of preamble and gap; frames shorter than 64 bytes are
padded; tail drop drops whole frames; a queue is FIFO; every random draw is seeded; the C++ queue reproduces the Python one to
the nanosecond.
\bfield{Acceptance tests} \texttt{code/firm/netsim/tests/}: frame sizes of known packets (64, 104, 1\,518 bytes); RFC~1112
address mapping and its 32-to-1 overlap; the peak-rate window on a hand-made stream; FIFO order, conservation of frames and
the buffer bound of the \omterm{def:nw:networking-for-trading:egress}{egress queue}; ten frames arriving together leave after ten serialisations; the fixture shared with the
C++ test.
\bfield{Stretch} Priority queues (strict priority for orders over market data on a shared uplink); a shared-buffer switch in
which ports borrow from a common pool under a dynamic threshold, to reproduce \cref{dat:nw:networking-for-trading:buffer}'s
arithmetic in simulation.
\end{build}

\begin{omsources}
\item OPRA, \emph{Key Operating Metrics of U.S. Options Securities Information Processor} (August 2026); SIAC, \emph{Revised
OPRA Capacity Projections} (15 September 2025).
\item G.~Fairhurst, ``Ethernet Frame Calculations'' (University of Aberdeen); UNH InterOperability Laboratory, \emph{10
Gigabit Ethernet} MAC tutorial.
\item RFC~768 (UDP), RFC~791 (IPv4), RFC~9293 (TCP), RFC~896 (Nagle), RFC~1122 (host requirements), RFC~1112 and RFC~5771
(multicast addressing), RFC~3376 (IGMPv3), RFC~4541 (snooping); Linux \texttt{tcp(7)}.
\item Cisco, \emph{Nexus 3548-X, 3524-X, 3548-XL and 3524-XL Switches Data Sheet}.
\end{omsources}

\section{Exercises}

\begin{exercise}[$\star$]\label{exo:nw:networking-for-trading:1}
What is the \omterm{def:nw:networking-for-trading:ser}{serialisation delay} of a 256-byte frame at \qty{25}{\giga\bit\per\second}, and how many 64-byte frames can a
\qty{10}{\giga\bit\per\second} link carry in one second?
\end{exercise}

\begin{exercise}[$\star$]\label{exo:nw:networking-for-trading:2}
To which Ethernet address is the \omterm{def:nw:networking-for-trading:group}{multicast group} 233.54.12.111 mapped? Give another group in 224.0.0.0/4 that shares it.
\end{exercise}

\begin{exercise}[$\star$]\label{exo:nw:networking-for-trading:3}
From \cref{dat:nw:networking-for-trading:opra}, what bandwidth must a firm taking both copies of the feed, with the
retransmission allowance, provision for the 10-millisecond peak, and how many \qty{25}{\giga\bit\per\second} links is that?
\end{exercise}

\begin{exercise}[$\star\star$]\label{exo:nw:networking-for-trading:4}
A gateway writes a new order and, forty microseconds later, a cancel of another order on the same TCP connection, without
\texttt{TCP\_NODELAY}. The venue acknowledges data only with its own messages or after its delayed-acknowledgement timer. What
can happen to the cancel, and why is the order of the two never at risk?
\end{exercise}

\begin{exercise}[$\star\star$]\label{exo:nw:networking-for-trading:5}
Four inputs send at \qty{6}{\giga\bit\per\second} each for 50 microseconds to a \qty{10}{\giga\bit\per\second} port with an empty
queue. How many bytes are queued at the end of the burst, and how long does the last byte wait?
\end{exercise}

\begin{exercise}[$\star\star$]\label{exo:nw:networking-for-trading:6}
In \cref{fig:nw:networking-for-trading:opra}, compute the ratio of the busiest millisecond's rate to the busiest second's rate
in January 2024 and in July 2026. What changed, and what does it mean for a link sized in 2024?
\end{exercise}

\begin{exercise}[$\star\star\star$]\label{exo:nw:networking-for-trading:7}
\emph{Coding.} With \texttt{nw\_wire.zero\_loss\_buffer}, compute the mean zero-loss buffer over twenty seeds when 10\% and when
50\% of the packets come with common events. How does the buffer scale with the share?
\end{exercise}

\begin{exercise}[$\star\star\star$]\label{exo:nw:networking-for-trading:8}
\emph{Find the flaw.} ``Our market-data port averages \qty{3}{\giga\bit\per\second} on a \qty{10}{\giga\bit\per\second} link and
its one-minute peak is \qty{4.2}{\giga\bit\per\second}. We have more than half the link spare; the packet loss our handler
reports must be in the server.''
\end{exercise}

\section{Problem: The Open That Did Not Fit}

\begin{problem}[{Weekend problem --- eight feeds, one port and a shared buffer}]\label{pb:nw:networking-for-trading:1}
A firm's market-data switch sends eight venue feeds to one strategy server over one \qty{10}{\giga\bit\per\second} port. Each
feed averages \qty{0.6}{\giga\bit\per\second}; at a market-wide event all eight burst together at \qty{2.5}{\giga\bit\per\second}
each for one millisecond. Frames average 200 bytes. The switch is the one of \cref{dat:nw:networking-for-trading:buffer}; during
such an event every port of the block is busy, so the port gets its sixteenth of the shared 6~MB.

\medskip\noindent\textbf{Part I --- Averages.}
\begin{enumerate}
\item What is the port's average load and utilisation?
\item How many frames a second does it carry on average (preamble and gap included)?
\item What is the offered rate during the event, and the utilisation?
\item Why does the port's one-minute counter never show the event?
\end{enumerate}

\medskip\noindent\textbf{Part II --- The queue.}
\begin{enumerate}[resume]
\item How many bytes are queued at the end of the burst, with an unlimited buffer?
\item How many frames is that?
\item How long does the last byte of the burst wait in the queue?
\item After the burst the feeds return to their averages. How long does the queue take to drain?
\end{enumerate}

\medskip\noindent\textbf{Part III --- The buffer.}
\begin{enumerate}[resume]
\item What buffer would lose nothing?
\item How much buffer does the port actually get during the event?
\item How many bytes and how many frames are dropped?
\item What share of the burst's frames is that?
\item Would the whole 6~MB block have been enough?
\end{enumerate}

\medskip\noindent\textbf{Part IV --- The verdict.}
\begin{enumerate}[resume]
\item State the \emph{named result}: the minimum buffer for zero loss, and the frames dropped at the port's share.
\item How long would the burst have to last for the whole 6~MB block to overflow?
\item Would a \qty{25}{\giga\bit\per\second} port queue at all during the event?
\item The firm takes both copies of each feed on separate ports of the same block. Does that help?
\item Which switch counter shows the loss, and which handler counter?
\item How do the handler's recovery paths (One Quant Book~13, chapter~18) cope with the loss?
\item In one sentence: why is the event that fills one port the event that fills all of them?
\end{enumerate}
\end{problem}

\section{Interview questions}

\begin{interviewq}[$\star$ \iqroles{developer}]\label{iq:nw:networking-for-trading:1}
What is \omterm{def:nw:networking-for-trading:ser}{serialisation delay}? Why does going from 10 to \qty{25}{\giga\bit\per\second} matter for an order that fits in one small
frame?
\end{interviewq}

\begin{interviewq}[$\star$ \iqroles{developer, trader}]\label{iq:nw:networking-for-trading:2}
Why is market data sent over UDP multicast and orders over TCP?
\end{interviewq}

\begin{interviewq}[$\star\star$ \iqroles{developer}]\label{iq:nw:networking-for-trading:3}
Explain \omterm{def:nw:networking-for-trading:nagle}{Nagle's algorithm} and the delayed acknowledgement. What do you set on an order-entry socket, and why?
\end{interviewq}

\begin{interviewq}[$\star\star$ \iqroles{developer}]\label{iq:nw:networking-for-trading:4}
The feed handler reports gaps on both lines at the open, but the server's card shows no drops. Where do you look?
\end{interviewq}

\begin{interviewq}[$\star\star$ \iqroles{developer}]\label{iq:nw:networking-for-trading:5}
What does \omterm{def:nw:networking-for-trading:snooping}{IGMP snooping} do, and what happens in a trading network without it?
\end{interviewq}

\begin{interviewq}[$\star\star\star$ \iqroles{developer, researcher}]\label{iq:nw:networking-for-trading:6}
How would you size the buffer of a switch port that receives several market-data feeds?
\end{interviewq}
